% !TeX program = xelatex
% 编译：latexmk -xelatex -interaction=nonstopmode -halt-on-error main.tex
\documentclass[aspectratio=169,9pt]{beamer}
\usepackage{idea-theme}

% ===== 只需先修改这里 =====
\title[研究 Idea]{研究 Idea 标题}
\subtitle{用一句话说明：为解决什么问题，提出什么方法}
\author{汇报人：姓名}
\institute{课题组 / 学院\quad ·\quad 导师：姓名}
\date{\today}

\begin{document}

% 01 封面：与内页相同的浅色画布、白色卡片、细边框。
\begin{frame}[plain]
  \begin{tikzpicture}[remember picture,overlay]
    \begin{scope}[shift={(current page.south west)},x=1cm,y=1cm]
      \draw[idea panel] (1,1.3) rectangle (15,7.7);
      \node[anchor=west,inner sep=0,text=IdeaPurple,font=\fontsize{8}{11}\selectfont]
        at (2,6.75) {研究构想};
      % 辅助色只作为小型色标，与方法图的配色保持一致。
      \foreach \xx/\tone in {12.55/IdeaBlue,12.88/IdeaPurple,13.21/IdeaAmber,13.54/IdeaTeal,13.87/IdeaCoral}{
        \fill[\tone!65!white,rounded corners=.7pt] (\xx,6.65) rectangle ++(.17,.17);
      }
      \node[anchor=north west,inner sep=0,text width=12cm,align=left] at (2,5.8) {%
        {\color{IdeaInk}\fontsize{27}{34}\selectfont\bfseries\inserttitle\par}
        \vspace{4mm}
        {\color{IdeaGray}\fontsize{10}{14}\selectfont\insertsubtitle\par}
      };
      \draw[IdeaLine,line width=.4pt] (2,2.85)--(14,2.85);
      \node[anchor=west,inner sep=0,text=IdeaGray,font=\fontsize{9}{12}\selectfont]
        at (2,2.12) {\insertauthor};
      \node[anchor=east,inner sep=0,text=IdeaGray,font=\fontsize{8}{11}\selectfont]
        at (14,2.12) {\insertdate};
    \end{scope}
  \end{tikzpicture}
\end{frame}

\IdeaSection{01}{MOTIVATION}{问题与假设}
% 02 两张等高卡片：左侧定义，右侧证据。
\begin{frame}{从一个明确的问题出发}
  \begin{columns}[T,onlytextwidth]
    \begin{column}{.485\textwidth}
      \IdeaCard[31mm]{01 / TASK}{任务定义}{%
        给定输入 $x$ 与条件 $c$，\par
        预测 $y$ 或生成候选 $x^\star$。\par
        \vspace{3mm}\Hint{研究对象：[填写]\quad 数据范围：[填写]}}
    \end{column}
    \begin{column}{.485\textwidth}
      \IdeaCard[31mm]{02 / GAP}{关键瓶颈}{%
        [描述最关键的失效情形]\par
        [说明它对研究目标的影响]\par
        \vspace{3mm}\Hint{证据来源：[作者，年份，图 / 表]}}
    \end{column}
  \end{columns}
  \vfill
  \Takeaway{[用一句话定义本研究需要解决的具体问题。]}
\end{frame}

% 03 核心假设：不是提前宣布有效。
\begin{frame}{一个可检验的核心假设}
  \begin{center}
  \begin{tikzpicture}[idea diagram]
    \node[idea io] (a) at (0,0) {已知瓶颈\\[2pt]\textbf{[信息缺失]}};
    \node[idea solid] (b) at (4.6,0) {提出干预\\[2pt]\textbf{[引入关键机制]}};
    \node[idea primary] (c) at (9.2,0) {预期变化\\[2pt]\textbf{[目标指标改善]}};
    \draw[idea arrow] (a)--node[idea note,above]{干预}(b);
    \draw[idea arrow] (b)--node[idea note,above]{检验}(c);
  \end{tikzpicture}
  \end{center}
  \vspace{3mm}
  \begin{columns}[T,onlytextwidth]
    \begin{column}{.47\textwidth}
      \Tag{成立的前提}
      \small [关键输入可获得；约束合理；评价可信]
    \end{column}
    \begin{column}{.47\textwidth}
      \Tag{否定该假设的证据}
      \small [移除机制后无差异；域外评价不改善]
    \end{column}
  \end{columns}
  \vfill
  \Takeaway{[引入什么机制，预期改善什么；在哪种条件下不成立。]}
\end{frame}

\IdeaSection{02}{METHOD}{方法框架}
% 04 总框架：白色阶段卡片 + 核心模块强调 + 独立训练反馈。
\begin{frame}{方法框架}
  \begin{center}
  \begin{tikzpicture}[idea diagram,y=1.4cm]
    \node[idea data] (x) at (0,0) {输入\\$x,c$};
    \foreach \xx/\num/\en/\tone in {3/A/ENCODE/IdeaBlue,7/B/INTERACT/IdeaPurple,11/C/PREDICT/IdeaTeal}{
      \draw[idea panel,draw=\tone!25,fill=\tone!3!white] (\xx-1.55,-1.3) rectangle (\xx+1.55,1.5);
      \node[idea note,anchor=west,text=\tone] at (\xx-1.3,1.15) {\num};
      \node[idea note,anchor=east,font=\fontsize{6}{8}\selectfont] at (\xx+1.3,1.15) {\en};
    }
    \node[font=\bfseries\small,text=IdeaInk] at (3,.65) {表征构建};
    \node[font=\bfseries\small,text=IdeaInk] at (7,.65) {条件交互};
    \node[font=\bfseries\small,text=IdeaInk] at (11,.65) {任务输出};
    \node[idea data] (a) at (3,-.15) {$z=f_\theta(x)$};
    \node[idea innovation] (b) at (7,-.15) {$h=g_\phi(z,c)$};
    \node[idea result] (c) at (11,-.15) {$\hat y=q_\psi(h)$};
    \node[idea note] at (3,-.9) {编码\quad /\quad 对齐};
    \node[idea note] at (7,-.9) {融合\quad /\quad 约束};
    \node[idea note] at (11,-.9) {预测\quad /\quad 评估};
    \draw[idea arrow,draw=IdeaBlue] (x.east)--++(.35,0)|-(a.west);
    \draw[idea arrow,draw=IdeaBlue] (a)--node[idea note,above]{$z$}(b);
    \draw[idea arrow,draw=IdeaTeal] (b.east)--node[idea note,above]{$h$}(c.west);
    \draw[idea feedback,draw=IdeaCoral] (11,-1.3)--(11,-1.9)--(3,-1.9)--(3,-1.3);
    \node[idea note,text=IdeaCoral!85!black,fill=IdeaCanvas,inner xsep=3mm] at (7,-1.9)
      {$\mathcal L=\mathcal L_{\mathrm{task}}+\lambda\mathcal L_{\mathrm{reg}}$};
    \node[idea note,font=\fontsize{7}{10}\selectfont] at (7,-2.15) {训练反馈\quad /\quad 任务监督与约束联合优化};
  \end{tikzpicture}
  \end{center}
\end{frame}

% 05 模块 A：大图 + 窄侧栏，可复制作为模块详情页。
\begin{frame}{表征构建}
  \begin{columns}[c,onlytextwidth]
    \begin{column}{.68\textwidth}
      \centering
      \begin{tikzpicture}[idea diagram,y=1.4cm]
        \node[idea data] (x1) at (0,1) {输入视图 1\\$x_1$};
        \node[idea data] (x2) at (0,-1) {输入视图 2\\$x_2$};
        \node[idea data] (e1) at (3,1) {编码器 $f_1$\\特征 $z_1$};
        \node[idea data] (e2) at (3,-1) {编码器 $f_2$\\特征 $z_2$};
        \node[idea data,fill=IdeaBlue!22!white] (z) at (6.2,0) {对齐 / 聚合\\[3pt]统一表征 $z$};
        \draw[idea arrow,draw=IdeaBlue] (x1)--(e1);
        \draw[idea arrow,draw=IdeaBlue] (x2)--(e2);
        \coordinate (z-upper) at ($(z.north west)!.3!(z.south west)$);
        \coordinate (z-lower) at ($(z.north west)!.7!(z.south west)$);
        \coordinate (encode-lane) at ($(e1.east)!.5!(z.west)$);
        \IdeaRouteX[draw=IdeaBlue]{e1.east}{z-upper}{encode-lane}
        \IdeaRouteX[draw=IdeaBlue]{e2.east}{z-lower}{encode-lane}
      \begin{scope}[on background layer]\node[idea dashed group,draw=IdeaBlue!55,fit=(x1)(x2)(e1)(e2)(z)] (encoder) {};
        \node[idea group label,text=IdeaBlue] at (encoder.north west) {多源表征模块};\end{scope}
      \end{tikzpicture}
    \end{column}
    \begin{column}{.28\textwidth}
      \Tag{设计说明}
      \small\textbf{输入}\par [来源、维度、预处理]\par\vspace{3mm}
      \textbf{关键选择}\par [为什么采用该编码方式]\par\vspace{3mm}
      \textbf{输出}\par $z\in\mathbb R^d$，供模块 B 使用
    \end{column}
  \end{columns}
\end{frame}

% 06 模块 B：两路输入 / 融合 / 残差结构。
\begin{frame}{条件感知的特征交互}
  \begin{center}
  \begin{tikzpicture}[idea diagram,y=1.4cm]
    \node[idea data] (z) at (0,0) {输入表征 $z$};
    \node[idea data] (p) at (3.3,0) {特征变换\\$u=Pz$};
    \node[idea innovation] (g) at (7.1,0) {条件调制\\$v=\gamma(c)\odot u$};
    \node[idea result] (h) at (10.7,0) {融合输出\\$h=u+v$};
    \node[idea condition] (cond) at (7.1,1.55) {条件编码 $\gamma(c)$};
    \draw[idea arrow,draw=IdeaBlue] (z)--(p);
    \draw[idea arrow] (p.east)--(g.west);
    \draw[idea arrow,draw=IdeaTeal] (g.east)--(h.west);
    \draw[idea arrow,draw=IdeaAmber] (cond.south)--(g.north);
    \draw[idea arrow,draw=IdeaBlue] (p.south)--(3.3,-1.15)--node[idea note,below]{保留原始特征的残差路径}(10.7,-1.15)--(h.south);
  \coordinate (lower) at (7,-1.65);
  \begin{scope}[on background layer]\node[idea panel,fit=(z)(p)(g)(h)(cond)(lower),inner sep=7pt] {};\end{scope}
      \end{tikzpicture}
  \end{center}
  % 架构示例：u,v,h,gamma(c) 同属 R^d；odot 为逐元素乘法。
\end{frame}

% 07 闭环结构：可用于主动学习 / 优化 / 迭代设计。
\begin{frame}{从预测到验证的闭环}
  \begin{columns}[c,onlytextwidth]
    \begin{column}{.66\textwidth}
      \centering
      \begin{tikzpicture}[idea diagram,y=1.4cm]
        \node[idea result] (a) at (0,1.05) {模型预测\\候选及不确定性};
        \node[idea condition] (b) at (4.5,1.05) {候选选择\\约束与采样策略};
        \node[idea validation] (c) at (4.5,-1.2) {独立验证\\实验 / 仿真};
        \node[idea data] (d) at (0,-1.2) {更新训练集\\重新训练};
        \draw[idea arrow,draw=IdeaTeal] (a)--(b);
        \draw[idea arrow,draw=IdeaAmber] (b)--(c);
        \draw[idea feedback,draw=IdeaCoral] (c)--(d);
        \draw[idea feedback,draw=IdeaBlue] (d)--(a);
      \begin{scope}[on background layer]\node[idea panel,fit=(a)(b)(c)(d),inner sep=7pt] {};\end{scope}
      \end{tikzpicture}
    \end{column}
    \begin{column}{.30\textwidth}
      \Tag{闭环边界}
      \small\textbf{选择规则}\par [目标与约束如何权衡]\par\vspace{3mm}
      \textbf{停止条件}\par [预算耗尽 / 达到阈值]\par\vspace{3mm}
      \textbf{评价边界}\par 独立测试集不参与反馈
    \end{column}
  \end{columns}
\end{frame}

% 08 训练与推理泳道图。
\begin{frame}{训练与推理}
  \begin{center}
  \begin{tikzpicture}[idea diagram,y=1.4cm]
    \node[idea note,font=\small\bfseries,text=IdeaPurple] at (-1.1,1) {训练};
    \node[idea note,font=\small\bfseries,text=IdeaPurple] at (-1.1,-1.4) {推理};
    \node[idea data] (tr) at (1,1) {训练输入及标签\\$(x,c,y)$};
    \node[idea primary] (mo) at (5.1,1) {前向预测 + 损失\\更新 $\theta,\phi,\psi$};
    \node[idea validation] (ck) at (9.3,1) {验证集选型\\冻结参数};
    \node[idea data] (te) at (1,-1.4) {新输入\\$(x_{\mathrm{new}},c_{\mathrm{new}})$};
    \node[idea primary] (fix) at (5.1,-1.4) {固定模型\\不使用测试标签};
    \node[idea result] (out) at (9.3,-1.4) {输出结果\\预测 / 候选};
    \draw[idea arrow,draw=IdeaBlue] (tr)--(mo);
    \draw[idea arrow,draw=IdeaCoral] (mo)--(ck);
    \draw[idea arrow,draw=IdeaBlue] (te)--(fix);
    \draw[idea arrow,draw=IdeaTeal] (fix)--(out);
    \draw[idea feedback,draw=IdeaCoral] (ck.south)--++(0,-.5)-|(fix.north);
  \begin{scope}[on background layer]\node[idea panel,fit=(tr)(mo)(ck)(te)(fix)(out),inner sep=7pt] {};\end{scope}
      \end{tikzpicture}
  \end{center}
\end{frame}

\IdeaSection{03}{EVIDENCE}{创新与验证}
% 09 创新点：不把组件罗列当作创新。
\begin{frame}{预期创新与证据对应}
  \renewcommand{\arraystretch}{1.55}
  \small
  \Surface{\begin{tabularx}{\linewidth}{@{}p{22mm}>{\raggedright\arraybackslash}X>{\raggedright\arraybackslash}X>{\raggedright\arraybackslash}X@{}}
    \toprule
    \textcolor{IdeaPurple}{设计层面} & \textcolor{IdeaPurple}{已有处理} & \textcolor{IdeaPurple}{提出的改变} & \textcolor{IdeaPurple}{验证方式}\\
    \midrule
    表征 / 模块 A & [基线表征] & [信息与归纳偏置] & 表征对照、等参数量\\
    交互 / 模块 B & [常规融合机制] & [具体算子或结构] & 交互消融与机制分析\\
    验证 / 模块 C & [现有验证协议] & [边界或闭环策略] & 独立测试、等预算\\
    \bottomrule
  \end{tabularx}}
  \vfill
  \Takeaway{[最重要的贡献是什么，哪些设计为其提供支撑。]}
\end{frame}

% 10 实验设计：按研究问题组织。
\begin{frame}{三个问题，三组实验}
  \renewcommand{\arraystretch}{1.5}
  \small
  \Surface{\begin{tabularx}{\linewidth}{@{}p{10mm}Xp{32mm}p{30mm}@{}}
    \toprule
    & \textcolor{IdeaPurple}{研究问题} & \textcolor{IdeaPurple}{对照设计} & \textcolor{IdeaPurple}{证据 / 指标}\\
    \midrule
    RQ1 & 整体方法是否有效？ & 代表基线、同预算 & 主指标与成本\\
    RQ2 & 核心模块是否必要？ & 完整模型与关键消融 & 指标变化、机制分析\\
    RQ3 & 在什么条件下失效？ & 域外 / 噪声 / 少样本 & 泛化差距、失败案例\\
    \bottomrule
  \end{tabularx}}
  \vspace{4mm}
  \Hint{数据：[样本量与划分]\quad 重复：[次数]\quad 预算：[填写]}
  \vfill
  \Takeaway{所有对照使用一致的数据边界和评价协议，避免泄漏与预算不公平。}
\end{frame}

% 11 结果表：不填入虚构数据；完成实验后替换“待填”。
\begin{frame}{主要实验结果}
  \begin{columns}[T,onlytextwidth]
    \begin{column}{.66\textwidth}
      \small\renewcommand{\arraystretch}{1.4}
      \Surface{\begin{tabularx}{\linewidth}{@{}Xccc@{}}
        \toprule
        方法 & 误差 $\downarrow$ & 成功率 $\uparrow$ & 成本 $\downarrow$\\
        \midrule
        基线 1 & \Pending & \Pending & \Pending\\
        基线 2 & \Pending & \Pending & \Pending\\
        \textcolor{IdeaPurple}{提出的方法} & \Pending & \Pending & \Pending\\
        \bottomrule
      \end{tabularx}}
      \vspace{1.5mm}
      \Hint{统计：[均值与波动]\quad 重复：[次数]}
    \end{column}
    \begin{column}{.29\textwidth}
      \Tag{结果解读}
      \small\textbf{观察}\par [实际提升或退化]\par\vspace{3mm}
      \textbf{解释}\par [证据支持的机制]\par\vspace{3mm}
      \textbf{边界}\par [尚不能支持的结论]
    \end{column}
  \end{columns}
  \vfill
  \Takeaway{[填写真实结果支持的结论，并保留必要的适用条件。]}
\end{frame}

% 12 结果图版式：原生 TikZ 示意曲线；无真实测量值。
\begin{frame}{性能变化来自哪里}
  \begin{columns}[c,onlytextwidth]
    \begin{column}{.56\textwidth}
      \centering
      \begin{tikzpicture}[idea diagram]
        \fill[white,rounded corners=3pt] (-.45,-.8) rectangle (6.6,3.85);
        \foreach \yy in {1,2,3} {\draw[IdeaLine!65,line width=.3pt] (0,\yy)--(6,\yy);}
        \draw[->,draw=IdeaGray] (0,0)--(6.25,0) node[below,idea note]{预算};
        \draw[->,draw=IdeaGray] (0,0)--(0,3.55) node[above,idea note]{性能};
        \draw[IdeaGray,dashed,line width=1pt] plot[smooth] coordinates {(.3,.4)(1.4,1)(2.6,1.35)(4,1.55)(5.7,1.6)};
        \draw[IdeaPurple,line width=1.4pt] plot[smooth] coordinates {(.3,.5)(1.4,1.3)(2.6,2)(4,2.5)(5.7,2.7)};
        \node[idea note,text=IdeaPurple,anchor=east] at (5.8,2.95) {完整模型（示意）};
        \node[idea note,anchor=east] at (5.8,1.25) {消融模型（示意）};
        \node[idea note] at (3,-.65) {仅展示图形样式 · 无真实数值};
      \end{tikzpicture}
    \end{column}
    \begin{column}{.39\textwidth}
      \Tag{消融表}
      \small\renewcommand{\arraystretch}{1.4}
      \Surface{\begin{tabularx}{\linewidth}{@{}Xc@{}}
        \toprule 变体 & 主指标\\\midrule
        完整方法 & \Pending\\
        移除模块 B & \Pending\\
        替换为简单融合 & \Pending\\
        移除约束项 & \Pending\\\bottomrule
      \end{tabularx}}
      \vspace{3mm}
      \Hint{控制：划分、预算、调参协议。}
    \end{column}
  \end{columns}
\end{frame}

\IdeaSection{04}{DISCUSSION}{讨论与决策}
% 13 讨论页：等高状态卡片，突出一个待决定的问题。
\begin{frame}{下一步，验证最关键的假设}
  \begin{columns}[T,onlytextwidth]
    \begin{column}{.315\textwidth}
      \IdeaCard[29mm]{01 / OBSERVED}{已有证据}{[调研或实验结论]\par\vspace{3mm}\Hint{依据：[记录 / 图表]}}
    \end{column}
    \begin{column}{.315\textwidth}
      \IdeaCard[29mm]{02 / UNKNOWN}{关键不确定性}{[核心假设的风险]\par\vspace{3mm}\Hint{影响：[结论 / 可行性]}}
    \end{column}
    \begin{column}{.315\textwidth}
      \IdeaCard[29mm]{03 / NEXT}{最小验证实验}{[最先执行的实验]\par\vspace{3mm}\Hint{周期：[填写]\quad 判据：[填写]}}
    \end{column}
  \end{columns}
  \vfill
  \Takeaway{待讨论：[继续推进该方向之前，最需要确认哪一项判断？]}
\end{frame}

% 14 完整框架示例：节点按内容自适应，分组由 fit 自动计算边界。
\IdeaSection{05}{TOOLKIT}{备用页}
\begin{frame}{完整框架：多源编码、检索增强与联合优化}
  \centering
  \begin{tikzpicture}[idea diagram,node distance=7mm,font=\fontsize{7.5}{9.5}\selectfont]
    \node[idea data] (x1) at (0,1.25) {主输入 $x_1$};
    \node[idea data,font=\scriptsize] (x2) at (0,0) {辅助输入 $x_2$};
    \node[idea condition] (cond) at (0,-1.45) {条件 / 先验 $c$};
    \node[idea feature] (e1) at (2.5,1.25) {主干编码器\\$z_1=f_1(x_1)$};
    \node[idea feature,font=\scriptsize] (e2) at (2.5,0) {轻量编码器\\$z_2=f_2(x_2)$};
    \node[idea memory,font=\scriptsize] (bank) at (2.5,-1.45) {外部知识库\\参考样本 / 原型};
    \node[idea innovation,font=\small] (fuse) at (5.5,.65) {条件交互\\$h=g(z_1,z_2,c,r)$};
    \node[idea memory] (retrieve) at (5.5,-1.45) {相似性检索\\上下文 $r$};
    \node[idea result] (head1) at (8.6,1.25) {预测头\\$\hat y=q(h)$};
    \node[idea result,font=\scriptsize] (head2) at (8.6,0) {不确定性估计\\$\sigma(h)$};
    \node[idea selection] (select) at (8.6,-1.45) {候选筛选\\约束 / 采样};
    \node[idea validation,font=\scriptsize] (loss) at (11.25,1.25) {联合损失\\$\mathcal L_{task}+\lambda\mathcal L_{reg}$};
    \node[idea validation] (verify) at (11.25,-1.45) {独立验证\\实验 / 仿真};
    % 左侧独立输入端口，避免角度锚点随长宽比变化。
    \coordinate (feature-upper) at ($(fuse.north west)!.28!(fuse.south west)$);
    \coordinate (feature-lower) at ($(fuse.north west)!.72!(fuse.south west)$);
    \coordinate (condition-port) at ($(fuse.south west)!.25!(fuse.south east)$);
    \coordinate (feature-lane) at ($(e1.east)!.5!(fuse.west)$);
    \coordinate (condition-lane) at (0,-.72);
    \coordinate (feedback-lane) at (0,-2.7);
    \coordinate (output-lane) at ($(fuse.east)!.5!(head1.west)$);
    \coordinate (fork) at (output-lane |- fuse.east);
    \draw[idea arrow,draw=IdeaBlue] (x1.east)--(e1.west);
    \draw[idea arrow,draw=IdeaBlue] (x2.east)--(e2.west);
    \IdeaRouteX[draw=IdeaCyan]{e1.east}{feature-upper}{feature-lane}
    \IdeaRouteX[draw=IdeaCyan]{e2.east}{feature-lower}{feature-lane}
    \IdeaRouteY[draw=IdeaAmber]{cond.north}{condition-port}{condition-lane}
    \draw[idea arrow,draw=IdeaIndigo] (bank.east)--(retrieve.west);
    % 一个共享分发干线，避免重复叠画；实心点表示真实连接。
    \draw[idea bus,draw=IdeaTeal] (fuse.east)--(fork);
    \draw[idea bus,draw=IdeaTeal] (output-lane |- head1.west)--(output-lane |- head2.west);
    \node[idea junction,fill=IdeaTeal] at (fork) {};
    \draw[idea arrow,draw=IdeaTeal] (output-lane |- head1.west)--(head1.west);
    \draw[idea arrow,draw=IdeaTeal] (output-lane |- head2.west)--(head2.west);
    \draw[idea arrow,draw=IdeaIndigo] (retrieve.north)--(fuse.south);
    \draw[idea arrow,draw=IdeaOlive] (head2.south)--(select.north);
    \draw[idea arrow,draw=IdeaCoral] (head1.east)--(loss.west);
    \draw[idea arrow,draw=IdeaOlive] (select.east)--(verify.west);
    \IdeaRouteY[idea feedback,draw=IdeaCoral]{verify.south}{bank.south}{feedback-lane}
    \node[idea edge label,below] at (7,-2.7) {验证反馈 / 更新参考样本};
    \begin{scope}[on background layer]
      \node[idea dashed group,draw=IdeaCyan!65,fit=(x1)(x2)(e1)(e2)] (enc) {};
      \node[idea group label,text=IdeaCyan!85!black] at (enc.north west) {多源编码};
      \node[idea dashed group,draw=IdeaIndigo!65,fit=(bank)(retrieve)] (mem) {};
      \node[idea group label,text=IdeaIndigo] at (mem.south east) [anchor=north east,yshift=-2pt] {检索增强子模块};
      \node[idea dashed group,draw=IdeaTeal!65,fit=(head1)(head2)(select)] (heads) {};
      \node[idea group label,text=IdeaTeal] at (heads.north west) {多任务决策};
    \end{scope}
  \end{tikzpicture}
  \vfill
  \begin{tikzpicture}[idea diagram]
    \foreach \xx/\yy/\tone/\label in {0/0/IdeaBlue/数据输入,2.6/0/IdeaCyan/特征编码,5.2/0/IdeaPurple/常规交互,7.8/0/IdeaAmber/条件约束,10.4/0/IdeaIndigo/记忆检索,0/-.38/IdeaOlive/筛选采样,2.6/-.38/IdeaTeal/任务输出,5.2/-.38/IdeaCoral/验证反馈,7.8/-.38/IdeaNovel/创新}{
      \fill[\tone,rounded corners=.6pt] (\xx,\yy) rectangle ++(.13,.13);
      \node[idea note,anchor=west,text=IdeaInk] at (\xx+.2,\yy+.065) {\label};
    }
  \end{tikzpicture}
\end{frame}

\end{document}
